% !TeX root = ../../Book3.tex
% 纲要标题：## 第1章　素理想、根理想与素谱
% 纲要位置：工作记录/计划纲要.md，第 1912 行。

\chapter{素理想、根理想与素谱}
\label{b3:ch:01}
\BookChapterNumber{b3:ch:01}

\begin{chapterabstract}
模去一个素理想得到整环，模去一个极大理想得到域；两种商环保留的代数信息不同。本章把全部素理想组成一个拓扑空间，以理想的根描述闭集，以素理想的包含关系辨认一般点、闭点和不可约分量。商环与局部化的谱对应进一步说明环的运算怎样改变这个空间，而幂零元的例子说明底拓扑空间不能恢复全部环结构。
\end{chapterabstract}

\begin{learninggoals}
\begin{itemize}
\item 用商环辨认素理想与极大理想，说明选择原理在一般存在性证明中的作用。
\item 证明根理想的素理想交表示，区分幂零元、幂零理想与约化性。
\item 在闭集、基本开集和理想之间往返计算，辨认一般点、闭点与不可约分量。
\item 写出环同态诱导的素谱映射，并区分主开集与一般局部化的像。
\end{itemize}
\end{learninggoals}

设 \(k\) 为域。在 \(k[x,y]\) 中，模去 \((x,y)\) 就把两个变量都置为零，得到域 \(k\)；模去 \((x)\) 则只把 \(x\) 置为零，留下多项式环 \(k[y]\)：
\[
k[x,y]/(x,y)\cong k,\qquad k[x,y]/(x)\cong k[y].
\]
前者对应原点，后者保留整条直线上的多项式。包含 \((x)\) 的极大理想仍给出这条直线的闭点，例如每个 \(a\in k\) 都给出 \((x,y-a)\)。只收集极大理想时，遗漏的是 \((x)\) 本身作为一个点；它记录整条直线共同满足的关系，而没有再指定 \(y\) 的取值。素谱把这个点也纳入其中，\cref{b3:prop:generic-and-closed-points}将证明它的闭包恰为整条 \(V(x)\)，因而称它为这条不可约直线的一般点。这里的“点”不只意味着给变量代入一组数。

\ProseParagraphBreak
整数环也有这种区别。\((2)\) 的商为域，而 \((0)\) 的商仍是整环 \(\mathbb Z\)。在 \(\operatorname{Spec}\mathbb Z\) 中，\((0)\) 的闭包是全谱，各个 \((p)\)（\(p\) 为素数）却是闭点；理想的包含关系说明这些点之间怎样联系。

\ProseParagraphBreak

同一空间还会忘记一些方程中的幂次。例如 \((6)\) 与 \((12)\) 在 \(\operatorname{Spec}\mathbb Z\) 中给出相同的闭集 \(\{(2),(3)\}\)。两者的根都是 \((6)\)，因此包含它们的素理想相同，尽管原来的两个理想不同。根理想将精确描述闭集所记录的信息；商环中的幂零元则说明取根时失去了什么。

\ProseParagraphBreak

本章使用第一卷的理想、商环、多项式和环局部化，也沿用其中的 Noether 升链条件。一般极大对象的存在性明确调用第一卷附录A的 Zorn 引理\footnote{佐恩（Max August Zorn，1906-1993），德国裔美国数学家，提出与选择公理等价的极大原则。}。此时不需要概形、层或同调代数；“点”暂时就是一个素理想，几何的直观只用来帮助解释已经给出的代数定义。

\ProseParagraphBreak
可参阅松村英之的《Commutative Ring Theory》\cite[\S 1, pp. 1--3; \S 4, pp. 24--25]{Matsumura1989CommutativeRingTheory}。该书 \S 4 的 Exercises 4.3--4.5、4.9--4.12 分别讨论局部化谱、商环谱、拟紧性及 Noether 素谱的不可约分解\cite[p. 29]{Matsumura1989CommutativeRingTheory}。该书的根记号与本卷略有不同，本卷以 \(\operatorname{Nil}(R)\) 表示幂零根，并将其与 Jacobson 根\footnote{雅各布森（Nathan Jacobson，1910-1999），美国数学家，研究环论和 Lie 代数，证明了稠密性定理。}区分。后者的模论作用在第三章讨论。

% !TeX root = ../../Book3.tex
% 纲要标题：### 1.1 素理想与极大理想
% 纲要位置：工作记录/计划纲要.md，第 1816 行。
% 纲要细目（写作范围）：
% * 素理想及商整环
% * 极大理想及剩余域
% * 素理想的存在与回避

\section{素理想与极大理想}
\label{b3:sec:0101}

在交换环中，模去一个理想以后，可能得到整环，也可能直接得到域。这两种情形分别由素理想和极大理想刻画。它们将成为本卷所研究空间的点，因此先把代数条件、商环及存在性说清楚。除另有说明外，本卷的环都交换且含幺，同态保持单位；允许零环，但素理想和极大理想必须是真理想。

\subsection{素理想与商整环}
\label{b3:subsec:0101:01}

\begin{definition}\label{b3:def:prime-maximal}
设 \(R\) 为交换含幺环。若 \(R\) 的真理想 \(\mathfrak p\) 满足
\[
ab\in\mathfrak p\ \Longrightarrow\ a\in\mathfrak p\text{ 或 }b\in\mathfrak p
\qquad(a,b\in R),
\]
则称 \(\mathfrak p\) 为\term[sulixiang]{素理想}[prime ideal]。若 \(R\) 的真理想 \(\mathfrak m\) 与 \(R\) 之间不存在其他理想，则称 \(\mathfrak m\) 为\term[jidalixiang]{极大理想}[maximal ideal]。
\end{definition}

\begin{proposition}\label{b3:prop:prime-maximal-quotient}
设 \(R\) 为交换含幺环，\(\mathfrak p,\mathfrak m\) 为 \(R\) 的真理想。\(\mathfrak p\) 为素理想当且仅当 \(R/\mathfrak p\) 为整环；\(\mathfrak m\) 为极大理想当且仅当 \(R/\mathfrak m\) 为域。特别地，极大理想必为素理想。
\end{proposition}
\begin{proof}
真理想保证商中的单位元不为零。商中 \((a+\mathfrak p)(b+\mathfrak p)=0\) 等价于 \(ab\in\mathfrak p\)，故第一个断言正是无零因子的条件。

若 \(\mathfrak m\) 极大，对 \(a\notin\mathfrak m\) 有 \(\mathfrak m+(a)=R\)，可写 \(1=u+ra\)、\(u\in\mathfrak m\)，所以 \(r+\mathfrak m\) 是 \(a+\mathfrak m\) 的逆。反之，若商为域，设 \(J\) 严格包含 \(\mathfrak m\)，取 \(a\in J\setminus\mathfrak m\)。将商中 \(a+\mathfrak m\) 的逆提升为 \(b\in R\)，则 \(ab-1\in\mathfrak m\subseteq J\)，而 \(ab\in J\)，相减得 \(1\in J\)。所以 \(J=R\)，证明极大性。域是整环，再由第一个断言即得最后结论。
\end{proof}

这重述了第一卷命题10.2.6的商环判据。实际检验时，寻找一个方便辨认的商环，常比逐一检查所有乘积更直接。例如 \(k[x,y]/(x)\cong k[y]\)，所以 \((x)\) 素而不极大；\(k[x,y]/(x,y)\cong k\)，所以 \((x,y)\) 极大。\(\mathbb Z\) 的素理想为 \((0)\) 和 \((p)\)、\(p\) 为素数，其中只有后者极大。

\ProseParagraphBreak
还可用理想的乘积表述素性：若 \(I J\subseteq\mathfrak p\)，则 \(I\subseteq\mathfrak p\) 或 \(J\subseteq\mathfrak p\)。否则分别取 \(a\in I\setminus\mathfrak p\)、\(b\in J\setminus\mathfrak p\)，有 \(ab\in\mathfrak p\)，违反素性。这个形式将在闭集并的计算中使用。

\subsection{极大理想与剩余域}
\label{b3:subsec:0101:02}

设 \(R\) 为交换含幺环，\(\mathfrak m\) 为 \(R\) 的极大理想。域 \(\kappa(\mathfrak m)=R/\mathfrak m\) 称为 \(\mathfrak m\) 处的\term[shengyu-yu]{剩余域}[residue field]。若 \(R\) 是域 \(k\) 上的非零代数，这个域含有 \(k\) 的像，却未必等于 \(k\)。例如 \(\mathbb R[t]\) 中 \((t^2+1)\) 极大，剩余域为 \(\mathbb C\)，而不存在实数 \(a\) 使这个理想等于 \((t-a)\)。第六章将说明有限型代数在代数闭域上的情形为何不同。
\symidx{\(\kappa(\mathfrak m)=R/\mathfrak m\)：极大理想处的剩余域}

\begin{theorem}\label{b3:thm:maximal-existence}
设 \(R\) 为交换含幺环。\(R\) 的任一真理想都包含于某个极大理想。因而非零交换含幺环至少有一个素理想；零环没有素理想。
\end{theorem}
\begin{proof}
固定真理想 \(I\)，在包含 \(I\) 的全部真理想中按包含排序。此集合非空。非空链的并仍为理想，且不含一：若一在并中，它已在某个链成员中，矛盾。空链由 \(I\) 提供上界。因此第一卷附录A定理A.1.4的 Zorn 引理给出极大元，它也是 \(R\) 的极大理想。取 \(I=0\) 得非零环的结论。零环没有真理想，故两个定义都不适用。
\end{proof}

一个常用后果是：\(a\in R\) 可逆当且仅当它不属于任何极大理想。若不可逆，\((a)\) 为真理想，用\cref{b3:thm:maximal-existence}即可找到包含它的极大理想；反向则因为包含单位的理想必为全环。注意这里并不要求 \(R\) 为 Noether 环。

\subsection{素理想的存在与回避}
\label{b3:subsec:0101:03}

\begin{lemma}[避开乘法集]\label{b3:lem:prime-disjoint-multiplicative}
设 \(R\) 为交换含幺环，\(S\subseteq R\) 含一且对乘法封闭，\(I\) 为 \(R\) 中与 \(S\) 不交的理想。则存在与 \(S\) 不交的素理想 \(\mathfrak p\supseteq I\)。
\end{lemma}
\begin{proof}
考虑包含 \(I\) 且与 \(S\) 不交的理想。链的并仍与 \(S\) 不交，空链有上界 \(I\)，所以 Zorn 引理给出极大者 \(\mathfrak p\)。它不含一，因而为真理想。若 \(a,b\notin\mathfrak p\)，极大性使 \(\mathfrak p+(a)\)、\(\mathfrak p+(b)\) 分别含有 \(s=u+ra\in S\)、\(t=v+qb\in S\)。若 \(ab\in\mathfrak p\)，展开乘积即得 \(st\in\mathfrak p\)，与 \(st\in S\) 矛盾。因此 \(\mathfrak p\) 素。
\end{proof}

这里使用选择原理的地方，是在避开乘法集的理想族中取极大元；它的极大性只相对于这个理想族而言，未必是原环的极大理想。对具体多项式环，常可直接写出所需的素理想；上述论证则保证任意交换环中相应对象存在。可参阅松村英之的《Commutative Ring Theory》\cite[Theorems 1.1--1.2, pp. 2--3]{Matsumura1989CommutativeRingTheory}。

\begin{lemma}[有限素理想回避]\label{b3:lem:finite-prime-avoidance}
设 \(R\) 为交换含幺环，\(I\) 是 \(R\) 的理想，\(\mathfrak p_1,\ldots,\mathfrak p_r\) 是 \(R\) 的素理想，其中 \(r\ge1\)。若 \(I\subseteq\bigcup_{i=1}^r\mathfrak p_i\)，则 \(I\subseteq\mathfrak p_i\) 对某个 \(i\) 成立。
\end{lemma}
\begin{proof}
删去被其他素理想包含的成员不改变并集，可设余下者两两不可比。若 \(I\) 不包含于任何一个，选 \(a_i\in I\setminus\mathfrak p_i\)。对 \(j\ne i\)，不可比性使我们可以选 \(b_{ij}\in\mathfrak p_j\setminus\mathfrak p_i\)。令
\[
c_i=a_i\prod_{j\ne i}b_{ij},\qquad c=\sum_i c_i.
\]
素性使 \(c_i\notin\mathfrak p_i\)，而 \(c_i\in I\) 且属于每个 \(\mathfrak p_j\)、\(j\ne i\)。对固定 \(i\)，和式中其他各项都在 \(\mathfrak p_i\) 内，所以 \(c\notin\mathfrak p_i\)。于是 \(c\in I\) 却不在所给并集中，矛盾。只有一个成员时把空乘积取为一，同一论证仍成立。
\end{proof}

这称为\term[sulixiang-huibi]{素理想回避引理}[prime avoidance lemma]的有限素理想版本。它常用来选择同时避开若干素理想的元素：若理想 \(I\) 不包含于任何一个 \(\mathfrak p_i\)，则 \(I\) 中至少有一个元素不属于它们的并，即
\[
I\nsubseteq\mathfrak p_i\quad(1\le i\le r)
\quad\Longrightarrow\quad
I\setminus\bigcup_{i=1}^r\mathfrak p_i\ne\varnothing.
\]

% !TeX root = ../../Book3.tex
% 纲要标题：### 1.2 根与约化
% 纲要位置：工作记录/计划纲要.md，第 1920 行。
% 纲要细目（写作范围）：
% * 幂零根和根理想；用有限混合乘积与无限变量例子区分逐元素幂零、统一指数和理想幂零
% * 根理想作为素理想的交；把不含指定元素转为避开其幂组成的乘法集
% * 约化环与零维例子；无限域乘积的论证对全部素理想成立

\section{根与约化}
\label{b3:sec:0102}

给定交换含幺环 \(R\) 的理想 \(I\)，一个元素 \(a\) 属于所有包含 \(I\) 的素理想，当且仅当它在商环 \(R/I\) 中的像是幂零元，也就是某次幂 \(a^n\) 落在 \(I\) 中。这一事实把素理想的包含关系与商环中的幂零现象联系起来。不过，这些素理想并不记录最小的幂次 \(n\)，也不能恢复原来的理想 \(I\)。将 \(I\) 换成它的根，保留的正是这组素理想包含信息。

\subsection{幂零根与根理想}
\label{b3:subsec:0102:01}

\begin{definition}\label{b3:def:radical-reduced}
设 \(R\) 为交换含幺环，\(I\) 为 \(R\) 的理想。定义 \(I\) 的\term[lixiang-de-gen]{根}[radical of an ideal]为
\[
\sqrt I=\{\,a\in R:\text{存在 }n\ge1\text{ 使 }a^n\in I\,\}.
\]
若 \(I=\sqrt I\)，则称 \(I\) 为\term[genlixiang]{根理想}[radical ideal]。\(\sqrt{(0)}\) 称为 \(R\) 的\term[miling-gen]{幂零根}[nilradical]，记为 \(\operatorname{Nil}(R)\)。若 \(\operatorname{Nil}(R)=0\)，则称 \(R\) 为\term[yuehua-huan]{约化环}[reduced ring]。
\end{definition}
\symidx{\(\sqrt I\)：理想的根}
\symidx{\(\operatorname{Nil}(R)\)：环的幂零根}

\(\sqrt I\) 确实是理想。若 \(a^m,b^n\in I\)，展开 \((a+b)^{m+n-1}\)，每项至少含 \(a^m\) 或 \(b^n\)，故该幂属于 \(I\)；负号和环元素倍乘也保持条件。还直接得到
\[
I\subseteq\sqrt I,\qquad \sqrt{\sqrt I}=\sqrt I,
\qquad \sqrt{I\cap J}=\sqrt{IJ}=\sqrt I\cap\sqrt J.
\]
最后一式中，若 \(a^m\in I,a^n\in J\)，则 \(a^{m+n}\in IJ\)；另一方向由 \(IJ\subseteq I\cap J\) 给出。对理想根迭代时，只需把两个存在的正指数相乘，便回到 \(I\)。

\ProseParagraphBreak
商环的幂零根为 \(\sqrt I/I\)，所以 \(R/I\) 约化当且仅当 \(I\) 是根理想。特别地，\(R_{\mathrm{red}}=R/\operatorname{Nil}(R)\) 约化，称为 \(R\) 的约化商。它把全部幂零元素送到零，而不要求剩下的非零元素都可逆。
\symidx{\(R_{\mathrm{red}}\)：环的约化商}

\begin{definition}\label{b3:def:nilpotent-ideal}
设 \(R\) 为交换含幺环，\(J\) 为 \(R\) 的理想。对正整数 \(N\)，\(J^N\) 是全部 \(N\) 个 \(J\) 中元素的乘积所生成的理想。若存在正整数 \(N\) 使 \(J^N=0\)，则称 \(J\) 为\term[mi-ling-lixiang]{幂零理想}[nilpotent ideal]；满足该式的最小正整数称为 \(J\) 的\term[miling-zhishu]{幂零指数}[nilpotency index]。
\end{definition}

理想幂零要求任意 \(N\) 个理想元素的混合乘积都为零，强于只对同一个元素取幂。对理想 \(J\)，下面三个条件逐项变强：
\[
\begin{aligned}
&\forall u\in J\ \exists n(u)\ge1, &&u^{n(u)}=0;\\
&\exists N\ge1\ \forall u\in J, &&u^N=0;\\
&\exists N\ge1, &&J^N=0.
\end{aligned}
\]
第一个条件在英文文献中称为 nil ideal，只要求每个元素各自幂零；第二个条件要求元素幂零有统一上界；第三个条件才是理想幂零。最后一个条件蕴含前两个，反向一般不成立。

\ProseParagraphBreak
例如，设 \(k\) 为域，在商环
\[
R=k[\varepsilon,\delta]/(\varepsilon^2,\delta^2)
\]
中仍用 \(\varepsilon,\delta\) 表示变量的像，令 \(J=(\varepsilon,\delta)\)。这个商环有 \(k\)-基 \(1,\varepsilon,\delta,\varepsilon\delta\)，所以混合乘积 \(\varepsilon\delta\ne0\)。因此
\[
J^2=(\varepsilon\delta)\ne0,\qquad J^3=0,
\]
后一式是因为任意三个生成元的乘积都会重复 \(\varepsilon\) 或 \(\delta\)。若 \(\operatorname{char}k=2\)，每个 \(u\in J\) 甚至都满足 \(u^2=0\)，而 \(J^2\) 仍不为零。这说明同一指数下，单个元素的幂与不同元素的混合乘积不能混为一谈。

\begin{proposition}\label{b3:prop:finite-nil-ideal}
设 \(R\) 为交换含幺环，\(J=(a_1,\ldots,a_m)\) 为 \(R\) 的理想，其中 \(m\ge0\)。若 \(a_i^{n_i}=0\)，且 \(n_i\ge1\)（\(1\le i\le m\)），则
\[
J^N=0,\qquad N=1+\sum_{i=1}^m(n_i-1).
\]
因此 Noether 环的幂零根是幂零理想。
\end{proposition}
\begin{proof}
\(m=0\) 时，\(J=0\)、\(N=1\)，结论成立。设 \(m\ge1\)。\(J^N\) 由全部 \(N\) 个生成元的乘积线性生成。若每个 \(a_i\) 出现次数都小于 \(n_i\)，总个数至多为 \(N-1\)，矛盾。所以每个乘积含某个 \(a_i^{n_i}\)，均为零。Noether 环的幂零根作为理想有限生成，每个生成元又幂零，故适用第一项。
\end{proof}

没有有限生成性时，一个理想即使每个元素都幂零，也未必是幂零理想。取域 \(k\)，考虑
\[
R=k[x_1,x_2,\ldots]/(x_1^2,x_2^2,\ldots).
\]
令 \(J=(\bar x_1,\bar x_2,\ldots)\)。每个 \(u\in J\) 只用有限多个变量，设有 \(r\) 个；任何 \(r+1\) 个正次单项式的乘积都会重复一个变量，所以 \(u^{r+1}=0\)。但 \(\bar x_1\cdots\bar x_N\ne0\) 对任意 \(N\) 成立，因为平方关系生成的理想由含某个平方因子的单项式线性生成，不能包含这个不含平方因子的单项式。又因商 \(R/J=k\) 约化，全部幂零元都在 \(J\) 中。因此 \(J=\operatorname{Nil}(R)\) 而 \(J^N\ne0\) 对全部 \(N\) 成立。

\ProseParagraphBreak
若 \(\operatorname{char}k=2\)，这个例子还有更强的性质：每个 \(u\in J\) 都满足 \(u^2=0\)。因为 \(u\) 是有限个正次单项式的线性组合，特征二时平方把有限和变成各项平方之和，而每个正次单项式的平方都为零。所以元素幂零即使有统一上界，也不能保证理想的某个幂为零。

\subsection{根理想的素理想交表示}
\label{b3:subsec:0102:02}

\begin{theorem}\label{b3:thm:radical-prime-intersection}
设 \(R\) 为交换含幺环。任意理想 \(I\subseteq R\) 满足
\begin{equation}\label{b3:eq:radical-prime-intersection}
\sqrt I=\bigcap_{\mathfrak p\supseteq I,\ \mathfrak p\text{ 素}}\mathfrak p.
\end{equation}
空交取为 \(R\)，因而公式也适用于 \(I=R\)。
\end{theorem}
\begin{proof}
若 \(a^n\in I\subseteq\mathfrak p\)，反复使用素性得到 \(a\in\mathfrak p\)，故左边包含于右边。

反过来，对 \(a\notin\sqrt I\)，我们要找一个包含 \(I\) 却不含 \(a\) 的素理想。对素理想 \(\mathfrak p\)，不含 \(a\) 等价于不含 \(a\) 的任何非负整数次幂：正次幂属于 \(\mathfrak p\) 会由素性推出 \(a\in\mathfrak p\)，而 \(1\) 本就不属于真理想。因此取乘法集 \(S=\{1,a,a^2,\ldots\}\)，便将不含 \(a\) 转化为与 \(S\) 不交。由 \(a\notin\sqrt I\) 可知 \(S\cap I=\varnothing\)：没有正次幂落在 \(I\) 中，且 \(I\ne R\)，所以 \(1\notin I\)。\cref{b3:lem:prime-disjoint-multiplicative}给出包含 \(I\) 并避开 \(S\) 的素理想 \(\mathfrak p\)，特别是 \(a\notin\mathfrak p\)。因此 \(a\) 不在右边的交中，证明反向包含。
\end{proof}

取 \(I=0\)，可知一个元素幂零当且仅当它属于全部素理想。于是映射
\[
R\longrightarrow\prod_{\mathfrak p\text{ 素}}R/\mathfrak p,
\qquad a\longmapsto(a+\mathfrak p)_{\mathfrak p}
\]
的核恰为 \(\operatorname{Nil}(R)\)。约化环因而可以嵌入整环的乘积，但本身未必是整环：\(k\times k\) 已给出最简单的反例。

\ProseParagraphBreak
在一般环上，不能把上述交中的素理想直接换成极大理想。以 \(R=k[t]_{(t)}\) 为例，其唯一极大理想为 \((t)R\)，但 \(R\) 是整环，幂零根为零。因此所有极大理想的交，即 Jacobson 根，可以严格大于幂零根。第二章将详细解释这个局部化环。

\subsection{约化环与零维例子}
\label{b3:subsec:0102:03}

\(k[\varepsilon]/(\varepsilon^2)\) 只有一个素理想 \((\varepsilon)\)：素理想必须包含幂零元 \(\varepsilon\)，商为域 \(k\)。它的约化商是 \(k\)。因此只记录有哪些素理想，尚不能区分这两个环；它们的区别在于前者存在非零平方零元。

\ProseParagraphBreak
\(\mathbb Z/12\mathbb Z\) 的素理想是 \((\bar2)\)、\((\bar3)\)，其交为 \((\bar6)\)，故幂零根平方为零。约化商为 \(\mathbb Z/6\mathbb Z\cong\mathbb F_2\times\mathbb F_3\)。后者没有非零幂零元，却仍有零因子，例如 \(\bar2\bar3=0\)。约化性只排除幂零元，不排除两个不同非零元素的乘积为零。

\ProseParagraphBreak
设 \(R\) 为非零交换含幺环。若 \(R\) 的每个素理想都是极大理想，则称 \(R\) 为零维环；这等价于素理想之间不存在严格包含。一般 Krull 维数将在第九章定义。\footnote{克鲁尔（Wolfgang Krull，1899-1971），德国数学家，研究交换环中的素理想、局部化与维数。}上述双数环及有限个域的乘积都零维。

\ProseParagraphBreak
零维性却不能控制理想的降链。令 \(k\) 为域，取无限乘积环 \(R=\prod_{i\ge1}k\)，其中的理想
\[
J_n=\{\,(a_i)\in R\mid a_1=\cdots=a_n=0\,\}
\qquad(n\ge1)
\]
满足 \(J_{n+1}\subsetneq J_n\)：第 \(n+1\) 个坐标为一、其余坐标为零的元素属于 \(J_n\)，却不属于 \(J_{n+1}\)。因此 \(R\) 不是 Artin 环。

\ProseParagraphBreak
另一方面，\(R\) 的每个素理想都极大。这里不把素理想限定为坐标投影的核；无限乘积中还可以有其他素理想，下面的论证对任意素理想都适用。对任意元素 \(a=(a_i)\)，在非零坐标取 \(b_i=a_i^{-1}\)，在零坐标取 \(b_i=0\)，便有 \(a=a^2b\)。模去任一素理想后，所得商是整环；其中任意非零的 \(\bar a\) 都可从 \(\bar a=\bar a^2\bar b\) 消去一份 \(\bar a\)，得到 \(1=\bar a\bar b\)，所以这个商是域。这样，\(R\) 是零维环，却有无限严格理想降链。第三章将讨论 Artin 环及其与零维性的关系。

% !TeX root = ../../Book3.tex
% 纲要标题：### 1.3 素谱的拓扑
% 纲要位置：工作记录/计划纲要.md，第 1828 行。
% 纲要细目（写作范围）：
% * Zariski 闭集与基本开集
% * 拟紧性（quasi-compactness，不附加 Hausdorff 条件）
% * 不可约闭集、一般点与闭点
% * Noether 素谱与有限不可约分量

\section{素谱的拓扑}
\label{b3:sec:0103}

把素理想作为点，还要指定哪些点属于同一个闭集。理想的包含关系给出一种自然规则：把包含同一理想的素理想收集起来。这个拓扑通常不是 Hausdorff 拓扑\footnote{豪斯多夫（Felix Hausdorff，1868-1942），德国数学家，建立一般拓扑中的分离公理等基本概念。}，甚至不同点之间可能有特殊化关系；这些现象正是需要保留的代数信息。

\subsection{Zariski 闭集与基本开集}
\label{b3:subsec:0103:01}

\begin{definition}\label{b3:def:spectrum-zariski}
设 \(R\) 为交换含幺环。\(R\) 的\term[sulixiang-pu]{素理想谱}[prime spectrum]是集合
\[
\operatorname{Spec}R=\{\,\mathfrak p\subset R:\mathfrak p\text{ 是素理想}\,\}.
\]
对 \(R\) 的理想 \(I\)，记 \(V(I)=\{\,\mathfrak p\in\operatorname{Spec}R\mid I\subseteq\mathfrak p\,\}\)。以这些 \(V(I)\) 为闭集得到的拓扑称为\term[Zariski-tuopu]{Zariski 拓扑}[Zariski topology]。对单个元素或有限个元素，简记 \(V(f)=V\bigl((f)\bigr)\)、\(V(f_1,\ldots,f_r)=V\bigl((f_1,\ldots,f_r)\bigr)\)。对 \(f\in R\)，记
\[
D(f)=\operatorname{Spec}R\setminus V(f)
=\{\,\mathfrak p:f\notin\mathfrak p\,\},
\]
称为\term[jiben-kaiji]{基本开集}[basic open subset]，也称主开集\termidx[zhukaiji]{主开集}。
\end{definition}
\symidx{\(\operatorname{Spec}R\)：交换环的素理想谱}
\symidx{\(V(I)\)：包含理想 \(I\) 的素理想组成的闭集}
\symidx{\(D(f)\)：素谱中的基本开集}

这里的规则确实定义拓扑，因为
\begin{equation}\label{b3:eq:zariski-closed-identities}
\begin{gathered}
V(0)=\operatorname{Spec}R,\qquad V(R)=\varnothing,
\\
\bigcap_\lambda V(I_\lambda)=V\left(\sum_\lambda I_\lambda\right),
\qquad V(I)\cup V(J)=V(IJ).
\end{gathered}
\end{equation}
前三式来自包含关系。最后一式是因为素理想包含 \(IJ\) 当且仅当它包含 \(I\) 或 \(J\)：若两者都不包含，分别从 \(I,J\) 中选取不属于该素理想的元素，它们的乘积属于 \(IJ\)，与素性矛盾。任意理想之和中的每个元素仅为有限和，所以这也说明无限交闭集的条件成立。

\ProseParagraphBreak
由根的交表示，还有
\begin{equation}\label{b3:eq:zariski-radical-order}
V(I)\subseteq V(J)\quad\Longleftrightarrow\quad
\sqrt J\subseteq\sqrt I.
\end{equation}
因此 \(V(I)=V(\sqrt I)\)，根理想与闭子集反向一一对应。另一方面，\(D(f)\cap D(g)=D(fg)\)，且
\(\operatorname{Spec}R\setminus V(I)=\bigcup_{f\in I}D(f)\)，故基本开集构成开集的一组基。特别地，\(D(f)=\varnothing\) 当且仅当 \(f\) 幂零，\(D(f)=\operatorname{Spec}R\) 当且仅当 \(f\) 是单位。

\subsection{拟紧性}
\label{b3:subsec:0103:02}

设 \(X\) 为拓扑空间。若 \(X\) 的每个开覆盖都有有限子覆盖，则称 \(X\) 为\term[nijin-kongjian]{拟紧空间}[quasi-compact space]。本书不在此词中附加 Hausdorff 条件。

\begin{proposition}\label{b3:prop:spec-quasicompact}
设 \(R\) 为交换含幺环，\(f\in R\)。则 \(\operatorname{Spec}R\) 及其基本开集 \(D(f)\) 都拟紧。
\end{proposition}
\begin{proof}
先将任一开覆盖细化为基本开覆盖。若 \(\operatorname{Spec}R=\bigcup_\lambda D(f_\lambda)\)，则不存在包含全部 \(f_\lambda\) 的素理想。由\cref{b3:thm:maximal-existence}，理想 \((f_\lambda)_\lambda\) 必为 \(R\)，所以一是其中有限个元素的组合。这有限个基本开集已经覆盖全谱，再取它们所属的原开集便得有限子覆盖。

对 \(D(f)\) 的任意子空间开覆盖，可先细化成由基本开集与 \(D(f)\) 的交组成的覆盖；交仍是基本开集，因为 \(D(f)\cap D(g)=D(fg)\)。因此只需考虑 \(D(f)\subseteq\bigcup_\lambda D(g_\lambda)\)。补集的包含关系是 \(V\bigl((g_\lambda)_\lambda\bigr)\subseteq V(f)\)。由\eqref{b3:eq:zariski-radical-order}，某个 \(f^N\) 属于 \((g_\lambda)_\lambda\)，于是属于其中有限个 \(g_\lambda\) 生成的理想。取基本开集又得到有限子覆盖。\(D(f)\) 为空时同样成立。
\end{proof}

\subsection{不可约闭集、一般点与闭点}
\label{b3:subsec:0103:03}

基本开覆盖可以由有限条代数关系控制。闭集还有另一种结构问题：它是否由几个更小的闭集拼成？若不能，下面的判据会把这个闭集辨认为某个素理想的闭包。

\begin{definition}\label{b3:def:irreducible-space}
设 \(X\) 为非空拓扑空间。若 \(X\) 不能写成两个真闭子集之并，则称 \(X\) 为\term[bukeyue-kongjian]{不可约空间}[irreducible space]。闭子集的不可约性用其子空间拓扑理解。若 \(X\) 不能分成两个非空、不交的开集，则称 \(X\) 为连通；不可约性比连通性强，因为在不可约空间中任意两个非空开集必须相交。
\end{definition}

\begin{proposition}\label{b3:prop:spec-irreducible}
设 \(R\) 为交换含幺环，\(I\) 为 \(R\) 的理想。若闭集 \(V(I)\subseteq\operatorname{Spec}R\) 非空，则 \(V(I)\) 不可约当且仅当 \(\sqrt I\) 为素理想。
\end{proposition}
\begin{proof}
若 \(\sqrt I=\mathfrak p\) 素，且 \(V(I)=\bigl(V(I)\cap V(J)\bigr)\cup\bigl(V(I)\cap V(K)\bigr)\)，则 \(V(\mathfrak p)\subseteq V(JK)\)。由\eqref{b3:eq:zariski-radical-order}，\(JK\subseteq\mathfrak p\)，故 \(J\subseteq\mathfrak p\) 或 \(K\subseteq\mathfrak p\)。相应的一个闭子集就是整个 \(V(I)\)。

反过来，若 \(\sqrt I\) 不素，选 \(a,b\notin\sqrt I\)、\(ab\in\sqrt I\)。有 \(V(I)\subseteq V(ab)=V(a)\cup V(b)\)，但 \(V(I)\) 不包含于任何一边，否则根的交表示会使 \(a\) 或 \(b\) 属于 \(\sqrt I\)。这给出两个真闭子集之并，违反不可约性。
\end{proof}

设 \(X\) 为拓扑空间，\(x\in X\)。若单点集 \(\{x\}\) 在 \(X\) 中为闭集，则称 \(x\) 为\term[bidian]{闭点}[closed point]。素谱中还常有闭包比单点大的点；这种闭包可以直接由理想的包含关系读出。

\begin{proposition}\label{b3:prop:generic-and-closed-points}
设 \(R\) 为交换含幺环，\(\mathfrak p\) 为 \(R\) 的素理想。它在 \(\operatorname{Spec}R\) 中所对应点的闭包为 \(V(\mathfrak p)\)。因此它是闭点当且仅当 \(\mathfrak p\) 极大；每个不可约闭子集恰有一个闭包为它的点。
\end{proposition}
\begin{proof}
闭集 \(V(I)\) 包含点 \(\mathfrak p\)，当且仅当 \(I\subseteq\mathfrak p\)，这又保证 \(V(\mathfrak p)\subseteq V(I)\)。所以 \(V(\mathfrak p)\) 是包含该点的最小闭集。它只有该点，当且仅当没有更大的真素理想；利用极大理想存在性，这等价于 \(\mathfrak p\) 极大。若闭集不可约，\cref{b3:prop:spec-irreducible}把它唯一写成 \(V(\mathfrak p)\)，根理想与闭集的一一对应保证此 \(\mathfrak p\) 唯一。
\end{proof}

设 \(X\) 为拓扑空间，\(Z\subseteq X\) 为不可约闭子集。若点 \(\eta\in Z\) 满足 \(\overline{\{\eta\}}=Z\)，则称 \(\eta\) 为 \(Z\) 的\term[yiban-dian]{一般点}[generic point]。一般点是相对于指定闭集而言的，并不与闭点互斥：闭点也是其单点闭子集的一般点。例如整环的素谱有一般点 \((0)\)；在 \(\operatorname{Spec}\mathbb Z\) 中，它的闭包是全谱，所有 \((p)\) 则是闭点。不同点具有不同闭包，所以素谱满足 \(T_0\) 分离性质，却一般不满足每点都闭，更不必是 Hausdorff 空间。

\ProseParagraphBreak
设 \(R\) 为交换含幺环，\(\mathfrak p,\mathfrak q\) 为 \(R\) 的素理想。若 \(\mathfrak p\subseteq\mathfrak q\)，则称 \(\mathfrak q\) 是 \(\mathfrak p\) 的\term[teshuhua]{特殊化}[specialization]，并称 \(\mathfrak p\) 是 \(\mathfrak q\) 的\term[yibanhua]{一般化}[generalization]。由闭包公式，这等价于 \(\mathfrak q\in\overline{\{\mathfrak p\}}\)。素理想 \(\mathfrak p\) 作为 \(V(\mathfrak p)\) 的一般点，记录这个不可约闭子集共同满足的代数关系；但这个闭集未必在全谱中极大。例如域 \(k\) 上的 \(\operatorname{Spec}k[x,y]\) 本身不可约，而 \((x)\) 是其真闭子集 \(V(x)\) 的一般点。

\ProseParagraphBreak
一个只有两个点的谱已经能显示这种拓扑与通常坐标点的差别。设 \(k\) 为域，令 \(R=k[t]_{(t)}\)。任一非零元素都可写为 \(t^n u\)，其中 \(n\ge0\)、\(u\in R\) 为单位；这是把分子中的 \(t\) 因子提取出来，而分子余下的因子及分母都不在 \((t)\) 中。若该元素属于素理想，则 \(n\ge1\)，并由素性得该理想含 \(t\)。另一方面，取值同态 \(f/g\mapsto\dfrac{f(0)}{g(0)}\) 将 \(R\) 满射到 \(k\)，核为 \((t)R\)，所以 \((t)R\) 极大。\(R\) 是整环，故 \((0)\) 也是素理想。这就得到\cref{b3:fig:two-point-specialization}中的全部两点。

\begin{figure}[htbp]
\centering
\begin{tikzcd}[column sep=huge]
\eta=(0) \arrow[r,"\text{特殊化}"] & s=(t)R
\end{tikzcd}
\caption{两点素谱的特殊化关系}
\label{b3:fig:two-point-specialization}
\end{figure}

图中的箭头表示 \(s\) 属于 \(\eta\) 的闭包，不表示距离或路径。具体地，
\[
\overline{\{\eta\}}=\{\eta,s\},\qquad
\overline{\{s\}}=\{s\}.
\]
所以三个开集恰为 \(\varnothing\)、\(\{\eta\}\) 和 \(\{\eta,s\}\)。其中 \(\{\eta\}=D(t)\) 非空，却不含全谱的唯一闭点 \(s\)。在一般素谱中，闭点的剩余域也可能比基域大；第六、七章将在代数闭基域的有限型情形把闭点解释为通常的坐标点。

\subsection{Noether 素谱与有限不可约分量}
\label{b3:subsec:0103:04}

不可约闭子集可能严格包含在另一个不可约闭子集中。例如上面的 \(V(x)\) 就包含在不可约的 \(\operatorname{Spec}k[x,y]\) 内。设 \(X\) 为拓扑空间，若不可约闭子集 \(Z\subseteq X\) 不严格包含于任何更大的不可约闭子集，则称 \(Z\) 为 \(X\) 的\term[bukeyue-fenliang]{不可约分量}[irreducible component]。

\ProseParagraphBreak
设 \(R\) 为交换含幺环，若素理想 \(\mathfrak p\) 不严格包含 \(R\) 的任何其他素理想，则称 \(\mathfrak p\) 为\term[jixiao-sulixiang]{极小素理想}[minimal prime ideal]。由\cref{b3:prop:spec-irreducible}与\eqref{b3:eq:zariski-radical-order}，\(V(\mathfrak p)\) 是全谱的不可约分量，当且仅当 \(\mathfrak p\) 是极小素理想：任一更大的不可约闭集都形如 \(V(\mathfrak q)\)，其包含关系恰对应 \(\mathfrak q\subseteq\mathfrak p\)。因此所有素理想都给出各自闭包的一般点，只有极小素理想给出全谱不可约分量的一般点。

\ProseParagraphBreak
不可约分量是否只有有限个，则是另一个问题。设 \(X\) 为拓扑空间，若 \(X\) 的每条闭集降链最终稳定，则称 \(X\) 为\term[Noether-tuopu-kongjian]{Noether 拓扑空间}[Noetherian topological space]。在素谱中，闭集降链对应根理想升链；环的 Noether 条件使这些升链稳定，因而保证素谱是 Noether 拓扑空间。

\begin{proposition}\label{b3:prop:noether-spectrum-components}
设 \(R\) 为交换含幺 Noether 环。则 \(\operatorname{Spec}R\) 是 Noether 拓扑空间，并且是有限多个不可约分量之并。它们恰为 \(V(\mathfrak p_1),\ldots,V(\mathfrak p_r)\)，其中 \(\mathfrak p_i\) 是 \(R\) 的全部极小素理想；每个素理想都包含其中至少一个。零环对应空分解。
\end{proposition}
\begin{proof}
若 \(R=0\)，素谱为空，结论成立。以下设 \(R\ne0\)。

闭集降链由\eqref{b3:eq:zariski-radical-order}对应根理想升链，Noether 升链条件使其稳定。若存在不能表示为有限个不可约闭集之并的闭集，在这些闭集中取一个关于包含关系的极小元 \(Y\)。降链条件保证极小元存在，否则从任一成员出发连续选择更小者，便得到无限严格降链。于是 \(Y\) 的每个真闭子集都有有限不可约分解。\(Y\) 非空且不可能不可约，否则它自己就是一项分解，故可写成两个真闭子集 \(Y_1,Y_2\) 的并。分别分解 \(Y_1,Y_2\)，再将所得有限个不可约闭集合并，就得到 \(Y\) 的分解，矛盾。

因此全谱有有限不可约闭分解。删去重复成员和被其他成员包含者，得到两两不包含的 \(V(\mathfrak p_1),\ldots,V(\mathfrak p_r)\)，其中各 \(\mathfrak p_i\) 由\cref{b3:prop:spec-irreducible}为素理想。若非空不可约闭集 \(Z\) 被闭集 \(F_1,\ldots,F_r\) 覆盖，则
\[
Z=\bigcup_{i=1}^r(Z\cap F_i).
\]
对 \(r\) 归纳可知某个 \(Z\cap F_i=Z\)：\(r=1\) 时显然；\(r>1\) 时，将第一项与余下各项的并看作 \(Z\) 的两个闭子集。不可约性使其中一边等于 \(Z\)；若是后者，再对 \(r-1\) 项使用归纳。因此任一不可约闭集都包含于某个 \(V(\mathfrak p_i)\)。若不可约闭集 \(Z\) 包含 \(V(\mathfrak p_i)\)，便有
\(V(\mathfrak p_i)\subseteq Z\subseteq V(\mathfrak p_j)\)
对某个 \(j\) 成立；两两不包含性迫使 \(j=i\)，故 \(Z=V(\mathfrak p_i)\)。这说明所列成员就是全部不可约分量，也就是全部极小素理想对应的闭集。

最后，对任一素理想 \(\mathfrak q\)，不可约闭集 \(V(\mathfrak q)\) 包含于某个 \(V(\mathfrak p_i)\)，即 \(\mathfrak p_i\subseteq\mathfrak q\)。
\end{proof}

不同不可约分量可以相交。例如域 \(k\) 上 \(k[x,y]/(xy)\) 的两个极小素理想为 \((\bar x)\) 和 \((\bar y)\)，因为包含 \(xy\) 的素理想必含 \(x\) 或 \(y\)，而 \((x)\)、\((y)\) 本身为互不包含的素理想。两条坐标轴分量 \(V(\bar x)\)、\(V(\bar y)\) 在 \((\bar x,\bar y)\) 处相交，所以存在多个不可约分量并不意味着空间不连通。

% !TeX root = ../../Book3.tex
% 纲要标题：### 1.4 环同态与素谱
% 纲要位置：工作记录/计划纲要.md，第 1834 行。
% 纲要细目（写作范围）：
% * 素理想收缩与连续映射
% * 商环对应闭子空间
% * 主开集 \(D(f)\) 对应 \(R_f\)；一般局部化的谱对应与乘法集不交的素理想子空间，未必是开子空间
% ---

\section{环同态与素谱}
\label{b3:sec:0104}

环同态将元素沿一个方向送出，素理想则沿收缩反向移动。商环与局部化是理解这个反向关系的两种基本操作：前者给出闭子空间，后者保留避开指定分母的点。关于元素 \(f\) 的各次幂的局部化称为单元素局部化，它对应基本开集 \(D(f)\)；一般乘法集给出的子空间却未必是开集。

\subsection{素理想收缩与谱上的连续映射}
\label{b3:subsec:0104:01}

\begin{samepage}
\begin{proposition}\label{b3:prop:spectrum-functor}
设 \(R,S\) 为交换含幺环，\(\varphi:R\rightarrow S\) 为保单位的环同态。它给出映射
\[
\varphi^*:\operatorname{Spec}S\longrightarrow\operatorname{Spec}R,
\qquad\mathfrak q\longmapsto\varphi^{-1}(\mathfrak q).
\]
它连续。对每个理想 \(I\subseteq R\) 和每个元素 \(f\in R\)，有
\begin{equation}\label{b3:eq:spectrum-pullback}
(\varphi^*)^{-1}\bigl(V(I)\bigr)=V(IS),\qquad
(\varphi^*)^{-1}\bigl(D(f)\bigr)=D\bigl(\varphi(f)\bigr),
\end{equation}
其中 \(IS\) 是由全部 \(\varphi(a)\)、\(a\in I\)，在 \(S\) 中生成的理想。恒等同态诱导恒等映射。若 \(A\) 也是交换含幺环，\(\psi:S\rightarrow A\) 为保单位的环同态，则 \((\psi\circ\varphi)^*=\varphi^*\circ\psi^*\)。
\end{proposition}
\end{samepage}
\begin{proof}
因 \(1\) 不在 \(\mathfrak q\) 中，收缩是真理想；若 \(ab\) 的像在 \(\mathfrak q\)，素性使 \(\varphi(a)\) 或 \(\varphi(b)\) 在其中，所以收缩为素理想。包含 \(I\subseteq\varphi^{-1}(\mathfrak q)\) 等价于 \(IS\subseteq\mathfrak q\)，给出第一式和连续性。对主理想取补集得到第二式。最后两项由原像的恒等及复合规则直接得到。
\end{proof}

极大理想的收缩未必极大。例如 \(\mathbb Z\hookrightarrow\mathbb Q\) 把域的唯一素理想 \((0)\) 收缩为 \(\mathbb Z\) 的零理想；后者不是极大理想。这也是素谱采用全部素理想而不只采用极大理想的一个理由。

\subsection{商环与闭子空间}
\label{b3:subsec:0104:02}

\begin{theorem}\label{b3:thm:spec-quotient}
设 \(R\) 为交换含幺环，\(I\) 为 \(R\) 的理想。商映射 \(R\rightarrow R/I\) 诱导同胚
\[
\operatorname{Spec}(R/I)\longrightarrow V(I)\subseteq\operatorname{Spec}R,
\qquad\overline{\mathfrak p}\longmapsto\mathfrak p.
\]
这里 \(\mathfrak p\) 是 \(\overline{\mathfrak p}\) 的原像，逆映射为 \(\mathfrak p\mapsto\mathfrak p/I\)。
\end{theorem}
\begin{proof}
理想的商对应和商整环判据给出所列互逆双射。对包含 \(I\) 的理想 \(J\)，\(R/I\) 的闭集 \(V(J/I)\) 对应为 \(V(J)\)，它已包含于 \(V(I)\)；另一方面，\(V(I)\) 中任一子空间闭集可写为 \(V(I)\cap V(K)=V(I+K)\)，也如此得到。故双射及其逆都连续。
\end{proof}

取 \(I=\operatorname{Nil}(R)\)，所有素理想都包含它，因而得到 \(\operatorname{Spec}R_{\mathrm{red}}\cong\operatorname{Spec}R\)。底拓扑空间看不见被取商消去的幂零元。这不表示两个环相同；后面讨论加厚点与纤维时，会反复用到这一区别。

\subsection{主开集与一般局部化的谱}
\label{b3:subsec:0104:03}

设 \(R\) 为交换含幺环，\(T\subseteq R\) 为含一的乘法集，允许 \(0\in T\)，此时局部化为零环。第一卷定理13.2.2构造了环 \(T^{-1}R\)，命题13.2.3给出零判据：\(a/t=0\) 当且仅当存在 \(u\in T\) 使 \(ua=0\)。下面把第一卷定理13.4.3的局部化素理想对应补上拓扑说明。

\begin{theorem}\label{b3:thm:spec-localization}
设 \(R\) 为交换含幺环，\(T\subseteq R\) 为含一的乘法集。局部化映射 \(\iota:R\rightarrow T^{-1}R\) 诱导同胚
\[
\operatorname{Spec}(T^{-1}R)\longrightarrow
X_T=\{\,\mathfrak p\in\operatorname{Spec}R:\mathfrak p\cap T=\varnothing\,\},
\]
右侧赋予子空间拓扑。对应的逆为 \(\mathfrak p\mapsto T^{-1}\mathfrak p\)。特别地，对每个 \(f\in R\)，令 \(R_f=\{1,f,f^2,\ldots\}^{-1}R\)，则 \(\operatorname{Spec}R_f\cong D(f)\)。
\end{theorem}
\begin{proof}
局部化中的素理想不能包含可逆分母的像，故其收缩与 \(T\) 不交。若 \(\mathfrak p\cap T=\varnothing\)，记 \(\overline T\) 为 \(T\) 在 \(R/\mathfrak p\) 中的像。由第一卷定理13.2.4的局部化泛性质，商映射诱导满同态
\[
\theta:T^{-1}R\longrightarrow\overline T^{-1}(R/\mathfrak p),
\qquad a/t\longmapsto\overline a/\overline t.
\]
满射性来自分子、分母都能取原像。局部化的零判据说明，\(\theta(a/t)=0\) 当且仅当某个 \(u\in T\) 使 \(ua\in\mathfrak p\)；此时 \(a/t=(ua)/(ut)\in T^{-1}\mathfrak p\)。反过来，\(\theta\) 把生成扩张理想的每个 \(b/1\)、\(b\in\mathfrak p\)，都送到零。因此 \(\ker\theta=T^{-1}\mathfrak p\)，由环的第一同构定理得到
\[
(T^{-1}R)/(T^{-1}\mathfrak p)\cong\overline T^{-1}(R/\mathfrak p).
\]
右端是非零整环，因为 \(R/\mathfrak p\) 是整环且 \(\overline T\) 不含零。所以 \(T^{-1}\mathfrak p\) 为素理想。前面的核计算还说明，若 \(a/1\in T^{-1}\mathfrak p\)，则某个 \(ua\in\mathfrak p\)，其中 \(u\in T\)；因 \(u\notin\mathfrak p\)，素性迫使 \(a\in\mathfrak p\)。因此扩张后的收缩恢复 \(\mathfrak p\)。

反过来，设 \(\mathfrak q\) 为 \(T^{-1}R\) 的理想，并令 \(J=\iota^{-1}(\mathfrak q)\)。若 \(a/t\in\mathfrak q\)，则 \(a/1=(t/1)(a/t)\in\mathfrak q\)，所以 \(a\in J\)，从而 \(a/t\in T^{-1}J\)。若 \(a\in J\)、\(t\in T\)，则 \(a/t=(1/t)(a/1)\in\mathfrak q\)。故 \(T^{-1}J=\mathfrak q\)，证明了收缩后扩张也恢复原理想。这验证了素理想之间的双射。

由\cref{b3:prop:spectrum-functor}，该映射连续。局部化中的基本开集 \(D(a/t)\) 在上述双射下对应 \(D(a)\cap X_T\)：素理想是否含 \(a/t\)，只取决于它是否含分子。这些正是子空间的一组开集基，故逆映射也连续。取 \(T=\{1,f,f^2,\ldots\}\) 时，素理想与 \(T\) 不交恰好等价于不含 \(f\)，得到主开集的断言。
\end{proof}

一般的 \(X_T=\bigcap_{t\in T}D(t)\) 未必开。取 \(R=\mathbb Z\)、\(T=\mathbb Z\setminus\{0\}\)，其像只有一般点 \((0)\)。任何包含它的基本开集 \(D(n)\)、\(n\ne0\)，都还包含一个 \((p)\)，其中素数 \(p\) 不整除 \(n\)；这样的素数存在，因为有限多个素数不能穷尽全部素数。因此单点 \(\bigl\{(0)\bigr\}\) 不是开集，尽管它与 \(\operatorname{Spec}\mathbb Q\) 同胚。

\ProseParagraphBreak
\cref{b3:tab:spectrum-operations}将上述操作并列。表中 \(R\) 为交换含幺环，\(I\) 为理想，\(f\in R\)，\(\mathfrak p\) 为素理想，\(T\) 为含一的乘法集；各行都经收缩把新环的谱视为原谱的子空间。在 \(\mathfrak p\) 处局部化保留包含于 \(\mathfrak p\) 的素理想，取商 \(R/\mathfrak p\) 则保留包含 \(\mathfrak p\) 的素理想，两个方向相反。

\begin{table}[htbp]
\centering
\caption{取商与局部化在素谱中保留的点}
\label{b3:tab:spectrum-operations}
\begin{tabularx}{\textwidth}{@{}l>{\raggedright\arraybackslash}p{.29\textwidth}>{\raggedright\arraybackslash}X@{}}
\toprule
环的操作 & 原环中保留的素理想 \(\mathfrak q\) & 在 \(\operatorname{Spec}R\) 中的像 \\
\midrule
\(R\rightarrow R/I\) & \(I\subseteq\mathfrak q\) & \(V(I)\)，闭子空间 \\
\(R\rightarrow R_f\) & \(f\notin\mathfrak q\) & \(D(f)\)，基本开集 \\
\(R\rightarrow R_{\mathfrak p}\) & \(\mathfrak q\subseteq\mathfrak p\) & \(\mathfrak p\) 及其一般化组成的子空间，未必开 \\
\(R\rightarrow T^{-1}R\) & \(\mathfrak q\cap T=\varnothing\) & \(\bigcap_{t\in T}D(t)\)，未必开 \\
\bottomrule
\end{tabularx}
\end{table}

第二章将把这些环上的分式构造推广到模。


\begin{chaptersummary}
素性与极大性分别使商环成为整环与域。避开乘法集的极大性论证不仅给出素理想存在性，还证明 \(\sqrt I\) 等于包含 \(I\) 的全部素理想之交。这个交精确刻画 \(R/I\) 中哪些元素幂零，却不能恢复原理想 \(I\)，也不记录元素或幂零理想的具体幂零指数。有限生成的逐元素幂零理想必为幂零理想；没有有限生成条件时，即使每个元素的平方都为零，理想的各次幂仍可能都非零。

\ProseParagraphBreak
Zariski 闭集对应根理想，基本开集对应元素不被素理想包含。全谱和每个基本开集都拟紧；每个素理想都是其闭包的一般点，极大理想恰好给出闭点，而极小素理想给出全谱的不可约分量。Noether 性使不可约分量有限，却不能从底拓扑空间反推出原环一定是 Noether 环。\(k[t]/(t^m)\) 对不同 \(m\ge1\) 总有相同的单点谱，商环中 \(\bar t\) 的幂零指数却不同。

\ProseParagraphBreak
环同态使素理想反向收缩。取商给闭子空间，单元素局部化给基本开集，一般局部化则给避开全部分母的子空间，它未必开。下一章把这种局部观察推广到模，并说明怎样从各点处的结果恢复全局结论。
\end{chaptersummary}

\BookExercises

\begin{exercise}\label{b3:ex:01-axes-primes}
设 \(k\) 为域，\(R=k[x,y]/(xy)\)。证明 \((\bar x)\)、\((\bar y)\) 素而不极大，\((\bar x,\bar y)\) 极大。证明每个素理想至少包含 \(\bar x,\bar y\) 中一个。
\end{exercise}
\begin{solution}
前两个商分别是 \(k[y],k[x]\)，为整环而非域；第三个商为 \(k\)，由\cref{b3:prop:prime-maximal-quotient}得到结论。又 \(\bar x\bar y=0\) 属于任何素理想，素性迫使至少一个因子在其中。因此这两个不可约分量已经覆盖全谱。
\end{solution}

\begin{exercise}\label{b3:ex:01-residue-field-four}
求 \(\mathbb F_2[t]/(t^2+t+1)\) 的元素与非零元素的逆，说明对应理想是极大理想，但剩余域不是原系数域。
\end{exercise}
\begin{solution}
多项式在零、一处都不为零，故这个二次多项式不可约。商以 \(1,\alpha\) 为基，有四个元素 \(0,1,\alpha,1+\alpha\)，其中 \(\alpha^2=\alpha+1\)。等式 \(\alpha(1+\alpha)=1\) 给出后二者互逆，一自逆，所以商为四元域。商域判据给出极大性，元素个数则说明它不是 \(\mathbb F_2\)。
\end{solution}

\begin{exercise}\label{b3:ex:01-radical-powers}
设 \(k\) 为域。在 \(k[x,y]\) 中证明 \(\sqrt{(x^2,y^3)}=(x,y)\)，并求商环的幂零根的最小正幂零指数。
\end{exercise}
\begin{solution}
两个变量都有幂在原理想中，故 \((x,y)\) 包含于根；反向因原理想包含于极大理想 \((x,y)\)，其根也包含于其中。商中的幂零根 \(N=(\bar x,\bar y)\)。任一四次单项式都含 \(x^2\) 或 \(y^3\)，故 \(N^4=0\)；但 \(\bar x\bar y^2\in N^3\) 非零，因为 \((x^2,y^3)\) 中每个单项式都含 \(x^2\) 或 \(y^3\)，故 \(N^3\ne0\)。最小正指数为四。
\end{solution}

\begin{exercise}\label{b3:ex:01-radical-sum}
设 \(R\) 为交换含幺环，\(I,J\) 为 \(R\) 的理想。证明 \(\sqrt{I+J}=\sqrt{\sqrt I+\sqrt J}\)。再取域 \(k\)，以 \(I=(y),J=(y-x^2)\subseteq k[x,y]\) 说明两个根理想之和未必为根理想。
\end{exercise}
\begin{solution}
取根保持理想的包含关系，所以 \(I+J\subseteq\sqrt I+\sqrt J\) 给出
\[
\sqrt{I+J}\subseteq\sqrt{\sqrt I+\sqrt J}.
\]
另一方面，由 \(I,J\subseteq I+J\) 得 \(\sqrt I,\sqrt J\subseteq\sqrt{I+J}\)，故 \(\sqrt I+\sqrt J\subseteq\sqrt{I+J}\)。再取根，并用 \(\sqrt{I+J}\) 已是根理想，得到
\[
\sqrt{\sqrt I+\sqrt J}\subseteq\sqrt{I+J}.
\]
两式给出所求等式；这里取根保持包含，而闭集运算 \(V\) 才反转包含。例中两个商都为 \(k[x]\)，所以 \(I,J\) 为素理想，特别是根理想；但其和为 \((y,x^2)\)，其中 \(x^2\) 属于它而 \(x\) 不属于它，故和不为根理想。
\end{solution}

\begin{exercise}\label{b3:ex:01-units-nil-quotient}
设 \(R\) 为交换含幺环，\(J\subseteq R\) 为理想，且 \(J\) 的每个元素都幂零。证明 \(a\in R\) 可逆当且仅当 \(a+J\) 在 \(R/J\) 中可逆；这里不要求 \(J\) 本身为幂零理想。
\end{exercise}
\begin{solution}
只有反向需要说明。若 \(ab-1=u\in J\)，选 \(N\) 使 \(u^N=0\)，则 \(1+u\) 的逆为 \(\sum_{j=0}^{N-1}(-u)^j\)。所以 \(b\sum_{j=0}^{N-1}(-u)^j\) 是 \(a\) 的逆。只用了这一个元素 \(u\) 的幂零性，不需要对全部 \(J\) 同时选一个指数。
\end{solution}

\begin{exercise}\label{b3:ex:01-integer-radical}
求 \(\mathbb Z/18\mathbb Z\) 的全部素理想、幂零根和约化商。它的谱有几个点，哪些点闭？
\end{exercise}
\begin{solution}
由商环对应，素理想来自 \(\mathbb Z\) 中包含十八的素理想，即 \((2),(3)\)。故幂零根为两像理想的交 \((\bar6)\)，包含 \(0,\bar6,\bar{12}\)，其平方为零。约化商为 \(\mathbb Z/6\mathbb Z\)，谱有两个点，都是极大理想对应的闭点。两个单点的补集也闭，因此这个有限谱是离散的。
\end{solution}

\begin{exercise}\label{b3:ex:01-product-spectrum}
设 \(A,B\) 为交换含幺环。证明 \(\operatorname{Spec}(A\times B)\) 是 \(\operatorname{Spec}A\) 与 \(\operatorname{Spec}B\) 的不交并，且这两个部分都是开闭子集。
\end{exercise}
\begin{solution}
令 \(e=(1,0)\)、\(f=(0,1)\)。因 \(ef=0,e+f=1\)，素理想恰好包含二者之一。若包含 \(f\)，它形如 \(\mathfrak p\times B\)，且商为 \(A/\mathfrak p\)，故 \(\mathfrak p\) 必素；另一类为 \(A\times\mathfrak q\)。两部分分别是 \(D(e)=V(f)\)、\(D(f)=V(e)\)，均开闭。各部分的拓扑由商环同胚\cref{b3:thm:spec-quotient}识别为相应因子的谱，所以这不仅是集合双射。
\end{solution}

\begin{exercise}\label{b3:ex:01-infinite-prime-avoidance}
设 \(k\) 为域，令 \(I=(x,y)\subseteq k[x,y]\)。对全部常数项为零的不可约多项式 \(h\)，考察素理想 \((h)\)。证明它们的并包含 \(I\)，但其中没有一个包含整个 \(I\)。
\end{exercise}
\begin{solution}
域上的多项式环为唯一分解整环，这是第一卷定理12.2.6逐次添加两个变量的结论。对 \(0\ne f\in I\)，不可约分解的常数项乘积为零，所以某个不可约因子 \(h\) 的常数项为零；于是 \(f\in(h)\)。零也在这些理想中。若 \(I\subseteq(h)\)，则 \(h\) 同时整除 \(x,y\)，只能为单位，矛盾。对 \(n\ge1\)，商环 \(k[x,y]/(x-y^n)\cong k[y]\) 是整环，故 \(x-y^n\) 为素元，特别是不可约元；它们已经给出无限多个不同的此类不可约元。这个例子说明有限素理想回避不能直接推广到无限并。
\end{solution}

\begin{exercise}\label{b3:ex:01-closed-integers}
在 \(\operatorname{Spec}\mathbb Z\) 中求 \(V(30)\)、\(D(6)\) 及点集 \(\bigl\{(2),(3)\bigr\}\) 的闭包。再证明任意无限多个不同的闭点组成的集合稠密。
\end{exercise}
\begin{solution}
\(V(30)=\bigl\{(2),(3),(5)\bigr\}\)，\(D(6)\) 包含一般点 \((0)\) 及除 \((2),(3)\) 外的全部闭点。所给两个点的集合本身等于 \(V(6)\)，所以已经闭。若一个闭集 \(V(n)\) 包含无限多个不同闭点，则整数 \(n\) 被无限多个素数整除，只能 \(n=0\)。因此包含这些点的唯一闭集是全谱，它们的闭包即为全谱。
\end{solution}

\begin{exercise}\label{b3:ex:01-affine-line-spectrum}
设 \(k\) 为代数闭域。描述 \(\operatorname{Spec}k[t]\) 的点及全部真闭集，并证明两个非空开集必相交。
\end{exercise}
\begin{solution}
一元多项式环的非零素理想由不可约多项式生成；代数闭性使这些多项式都是一次式。因此点为 \((0)\) 及 \((t-a)\)、\(a\in k\)。任何真闭集形如 \(V(f)\)、\(f\ne0\)，恰为这个多项式的有限多个根对应的闭点。反过来，任意有限闭点集 \(\{(t-a_1),\ldots,(t-a_r)\}\) 都等于
\[
V\!\left(\prod_{i=1}^{r}(t-a_i)\right).
\]
空集对应空乘积 \(1\)，即 \(V(1)=\varnothing\)。一般点 \((0)\) 不属于任何真闭集，所以非空开集必含 \((0)\)，任意两个这样的开集因而相交。其不可约性也可由整环的零理想为素理想得到。
\end{solution}

\begin{exercise}\label{b3:ex:01-local-open-no-closed-point}
设 \(k\) 为域，\(R=k[t]_{(t)}\)。分别计算 \(R/(t^2)\) 与 \(R[1/t]\) 的素谱，以及它们在 \(\operatorname{Spec}R\) 中的像。判断这两个像是否包含 \(\operatorname{Spec}R\) 的闭点，说明它们虽同为单点空间，为何处于原谱的不同位置。
\end{exercise}
\begin{solution}
\(\operatorname{Spec}R\) 的两个点为 \(\eta=(0)\)、\(s=(t)R\)，其中只有 \(s\) 为闭点。由\cref{b3:thm:spec-quotient}，\(\operatorname{Spec}(R/(t^2))\) 的像是 \(V(t^2)=V(t)=\{s\}\)，所以它含有原谱的闭点。

由\cref{b3:thm:spec-localization}，\(\operatorname{Spec}(R[1/t])\) 的像为 \(D(t)=\{\eta\}\)，不含原谱的闭点。事实上，\(R\) 中每个非零元素都可写成 \(t^n u\)，其中 \(n\ge0\)、\(u\) 为单位；再使 \(t\) 可逆便使全部非零元素可逆，得到 \(R[1/t]=k(t)\)。两个源谱都是单点空间，但商映射保留的是闭点，局部化保留的是一般点。\(\eta\) 在开子空间 \(\{\eta\}\) 中是闭点，在原谱中却不是；闭点的含义取决于所在空间。
\end{solution}

\begin{exercise}\label{b3:ex:01-nonquasicompact-open}
设 \(k\) 为域，在 \(R=k[x_1,x_2,\ldots]\) 中令 \(U=\bigcup_{i\ge1}D(x_i)\)。证明 \(U\) 不拟紧，尽管全谱及每个 \(D(x_i)\) 都拟紧。
\end{exercise}
\begin{solution}
若有限个 \(D(x_i)\) 覆盖 \(U\)，取未在这有限指标中出现的 \(j\)。把所选变量送到零、\(x_j\) 送到一、其余变量送到零，得到满同态 \(R\rightarrow k\)。它的核为极大理想，属于 \(D(x_j)\subseteq U\)，却不在有限个所选基本开集内，矛盾。所以这个开覆盖没有有限子覆盖。
\end{solution}

\begin{exercise}\label{b3:ex:01-noether-topology-not-ring}
设 \(k\) 为域，\(V\) 为无限维 \(k\) 线性空间，在 \(R=k\oplus V\) 上规定 \((a,v)(b,w)=(ab,aw+bv)\)。证明 \(R\) 的谱只有一个点，但 \(R\) 不是 Noether 环。
\end{exercise}
\begin{solution}
对第三个元素 \((c,u)\)，乘法的两种结合次序都得到 \((abc,ab u+ac w+bc v)\)，单位为 \((1,0)\)。理想 \(0\oplus V\) 平方为零，商为 \(k\)，所以它是唯一素理想。谱因而是单点 Noether 空间。选无限线性无关序列 \(v_i\)，其有限初段所张成的 \(0\oplus\operatorname{span}_k\{v_1,\ldots,v_n\}\) 都是理想，并组成严格升链。故 \(R\) 不是 Noether 环；底拓扑空间无法检测全部环的有限性。
\end{solution}

\begin{exercise}\label{b3:ex:01-components-plane-line}
设 \(k\) 为域。在 \(k[x,y,z]\) 中证明 \((xy,xz)=(x)\cap(y,z)\)，并求 \(V(xy,xz)\) 的全部不可约分量及其一般点。
\end{exercise}
\begin{solution}
左边显然包含于右边。若 \(xh\in(y,z)\)，模去 \((y,z)\) 得 \(x\cdot h(x,0,0)=0\) 于整环 \(k[x]\)，所以 \(h(x,0,0)=0\)，即 \(h\in(y,z)\)。故反向包含成立。两个理想的商分别为 \(k[y,z]\) 与 \(k[x]\)，所以都素，且互不包含。因此不可约分量为 \(V(x)\)、\(V(y,z)\)，一般点分别为 \((x)\)、\((y,z)\)。它们的交为 \(V(x,y,z)\)，非空；不可约分量不是必然互不相交的部分。
\end{solution}

\begin{exercise}\label{b3:ex:01-thickened-point}
设 \(k\) 为域，\(m\ge1\) 为整数。比较环 \(k[t]/(t^m)\) 的素谱、\(k\) 维数和幂零根的幂零指数，说明谱同胚不能决定环同构。
\end{exercise}
\begin{solution}
每个素理想都含 \(\bar t\)，商为域，所以谱总为一个点。环以 \(1,\bar t,\ldots,\bar t^{m-1}\) 为基，维数为 \(m\)；\((\bar t)^m=0\)，若 \(m>1\)，其 \(m-1\) 次幂非零，因此最小正指数为 \(m\)。不同 \(m\) 的这些 \(k\) 代数因维数或幂零指数不同而不同构，但底谱相同。
\end{solution}

\begin{exercise}\label{b3:ex:01-maps-direction}
在交换含幺环及保单位环同态中，分别给出反例，说明环同态单射未必使谱映射单射，以及谱映射同胚未必使环同态单射或满射。
\end{exercise}
\begin{solution}
取域 \(k\)。包含 \(k\hookrightarrow k[t]\) 单射，谱上的映射却把多个点全部送到 \(\operatorname{Spec}k\) 的唯一点。商映射 \(k[\varepsilon]/(\varepsilon^2)\rightarrow k\) 不是单射，却在谱上给同胚。一个真域扩张 \(k\hookrightarrow K\) 不是满射，但两端的谱都是单点，也给同胚。环的元素及剩余域信息并不能由这个底空间映射单独恢复。
\end{solution}

\begin{exercise}\label{b3:ex:01-two-localizations}
描述 \(\operatorname{Spec}\mathbb Z[1/6]\) 和 \(\operatorname{Spec}\mathbb Z_{(2)}\) 在 \(\operatorname{Spec}\mathbb Z\) 中的像，判断两者是否开。
\end{exercise}
\begin{solution}
前者的像是 \(D(6)\)，为开集。后者只保留包含于 \((2)\) 的素理想，即 \(\bigl\{(0),(2)\bigr\}\)。任何包含一般点的开集都包含某个 \(D(n)\)、\(n\ne0\)，而这个基本开集还含无限多个其他素数点。因此后者不是开集。两种像均按\cref{b3:thm:spec-localization}具有正确的子空间拓扑。
\end{solution}

\begin{exercise}\label{b3:ex:01-square-map}
设 \(k\) 为域，\(b\in k\)，令 \(\varphi:k[s]\rightarrow k[t]\) 为满足 \(\varphi(s)=t^2\) 的 \(k\)-代数同态。计算闭点 \((s-b)\) 在 \(\varphi^*\) 下的原像及其补集。分别在 \(b\) 不是 \(k\) 中的平方、\(b=0\)、\(b=a^2\ne0\) 的情形下，求原像的点数及这些点的剩余域，并判断 \(k[t]/(t^2-b)\) 是否约化；最后一种情形须区分 \(\operatorname{char}k=2\) 与 \(\operatorname{char}k\ne2\)。
\end{exercise}
\begin{solution}
\((s-b)\) 是极大理想，所以 \(V(s-b)=\{(s-b)\}\)。由\eqref{b3:eq:spectrum-pullback}，所求原像及其补集分别为 \(V(t^2-b)\) 与 \(D(t^2-b)\)。包含 \(t^2-b\) 的素理想收缩后包含 \((s-b)\)，而收缩仍是真理想，故恰好收缩为这个闭点。

若 \(b\) 不是平方，则二次多项式 \(t^2-b\) 没有 \(k\) 中的根，因而不可约。原像只有 \((t^2-b)\) 一个点，剩余域是二次扩张 \(k[t]/(t^2-b)\)；这个商是域，所以约化。

若 \(b=0\)，原像只有 \((t)\)，剩余域为 \(k\)，但商 \(k[t]/(t^2)\) 中 \(\bar t\ne0\) 且 \(\bar t^2=0\)，因此不约化。若 \(b=a^2\ne0\) 且特征不为二，则
\[
t^2-b=(t-a)(t+a),\qquad a\ne-a.
\]
原像有 \((t-a)\)、\((t+a)\) 两个点，剩余域都为 \(k\)。两个因子互素，中国剩余定理给出商环 \(k\times k\)，所以约化。特征为二时，\(t^2-b=(t-a)^2\)，原像只有 \((t-a)\)，剩余域为 \(k\)，商中 \(t-a\) 的像非零而平方为零，所以不约化。原像同为单点时，剩余域和商环的幂零元仍可能不同。
\end{solution}

\begin{exercise}\label{b3:ex:01-reduction-universal}
设 \(R,S\) 为交换含幺环，且 \(S\) 约化。证明每个保单位环同态 \(R\rightarrow S\)，都唯一经 \(R_{\mathrm{red}}\) 分解。说明为何仅有谱同胚不能反向推出这个同态的泛性质。
\end{exercise}
\begin{solution}
幂零元的像仍幂零，而 \(S\) 约化，所以同态零化 \(\operatorname{Nil}(R)\)。定义 \(r+\operatorname{Nil}(R)\mapsto\varphi(r)\) 即得到良定义的环同态；原商映射满射，故分解唯一。这个泛性质谈的是到每个约化环的全部同态。例如 \(\mathbb Q\hookrightarrow\mathbb Q(i)\) 的谱同样是单点间同胚，若恒等映射 \(\mathbb Q\rightarrow\mathbb Q\) 能经它分解，就须有保单位环同态 \(\psi\) 使
\[
\mathbb Q\hookrightarrow\mathbb Q(i)\xrightarrow{\psi}\mathbb Q
\]
的复合为恒等映射。第一支箭头已经存在，第二支却不存在：由 \(i^2=-1\) 将得到 \(\psi(i)^2=-1\)，而有理数中没有这样的元素。因此底拓扑同胚没有提供这一代数泛性质。
\end{solution}

\begin{exercise}\label{b3:ex:01-irreducible-not-domain}
设 \(R\) 为交换含幺环。证明 \(\operatorname{Spec}R\) 不可约当且仅当 \(\operatorname{Nil}(R)\) 为素理想。给出一个非整环但谱不可约的非零环，并说明加入约化条件后能得到什么。
\end{exercise}
\begin{solution}
零环的谱为空，不可约性与幂零根为素理想这两项均不成立。对非零环，结论是\cref{b3:prop:spec-irreducible}在 \(I=0\) 时的应用。取域 \(k\)，双数环 \(k[\varepsilon]/(\varepsilon^2)\) 的幂零根 \((\varepsilon)\) 为极大理想，故谱不可约，但 \(\varepsilon\ne0\) 且平方为零，环不是整环。若另有约化性，幂零根为零，所以零理想素，商整环判据说明 \(R\) 本身是整环。
\end{solution}

\begin{exercise}\label{b3:ex:01-basic-open-certificate}
设 \(R\) 为交换含幺环，\(f\in R\)，\((g_\lambda)_{\lambda\in\Lambda}\) 为 \(R\) 中的一族元素。证明
\[
D(f)\subseteq\bigcup_{\lambda\in\Lambda}D(g_\lambda)
\]
当且仅当存在有限子集 \(F\subseteq\Lambda\)、整数 \(N\ge1\) 和元素 \(a_\lambda\in R\)（\(\lambda\in F\)），使
\[
f^N=\sum_{\lambda\in F}a_\lambda g_\lambda.
\]
再在 \(R=\mathbb Z\) 中取 \(f=6\)、\(g_1=8\)、\(g_2=20\)，求这样的最小 \(N\)，并写出一组系数。
\end{exercise}
\begin{solution}
令 \(J=(g_\lambda)_{\lambda\in\Lambda}\)。对所给开集包含关系取补集，得到 \(V(J)\subseteq V(f)\)。由\eqref{b3:eq:zariski-radical-order}，这等价于 \(f\in\sqrt J\)，即某个正次幂 \(f^N\) 属于 \(J\)。理想生成中的每个表达式只含有限项，因而恰好得到题设的有限等式。反过来，若该等式成立，任何包含全部 \(g_\lambda\) 的素理想都包含 \(f^N\)，从而包含 \(f\)，所以 \(V(J)\subseteq V(f)\)，再取补集便得开集包含关系。\(F\) 为空时，等式为 \(f^N=0\)，此时 \(D(f)=\varnothing\)。

整数例中 \((8,20)=(4)\)。\(6\) 不是四的倍数，而 \(6^2=36\) 是，所以最小指数为 \(N=2\)。等式
\[
36=2\cdot8+20
\]
给出系数 \(a_1=2\)、\(a_2=1\)，也给出无需逐个检查素理想的包含证书。
\end{solution}

\begin{exercise}\label{b3:ex:01-spectrum-image-closure}
设 \(R,S\) 为交换含幺环，\(\varphi:R\rightarrow S\) 为保单位的环同态，\(\varphi^*:\operatorname{Spec}S\rightarrow\operatorname{Spec}R\) 为收缩映射。证明
\[
\overline{\varphi^*(\operatorname{Spec}S)}=V(\ker\varphi).
\]
\end{exercise}
\begin{solution}
对 \(X=\operatorname{Spec}R\) 的任意点集 \(E\)，包含 \(E\) 的闭集 \(V(I)\) 恰好满足 \(I\subseteq\bigcap_{\mathfrak p\in E}\mathfrak p\)。因此
\[
\overline E=V\!\left(\bigcap_{\mathfrak p\in E}\mathfrak p\right),
\]
其中空交取为 \(R\)，也给出 \(\overline\varnothing=V(R)=\varnothing\)。取 \(E=\varphi^*(\operatorname{Spec}S)\)，由\cref{b3:thm:radical-prime-intersection}，
\[
\bigcap_{\mathfrak p\in E}\mathfrak p
=\bigcap_{\mathfrak q\in\operatorname{Spec}S}\varphi^{-1}(\mathfrak q)
=\varphi^{-1}\bigl(\operatorname{Nil}(S)\bigr).
\]
一个元素 \(a\in R\) 属于最后的理想，当且仅当某个 \(\varphi(a)^n=0\)，也就是 \(a^n\in\ker\varphi\)。故该理想等于 \(\sqrt{\ker\varphi}\)，而取根不改变零集，得到所需等式。
\end{solution}
